\documentclass[11pt, letterpaper]{article}

% --- PREÁMBULO ---
\usepackage[utf8]{inputenc}
\usepackage[T1]{fontenc}
\usepackage[letterpaper, top=2.2cm, bottom=2.2cm, left=2.2cm, right=2.2cm]{geometry}
\usepackage{xcolor}
\usepackage{graphicx}
\usepackage[most]{tcolorbox}
\usepackage{enumitem}
\usepackage{tikz}
\usepackage{booktabs}
\usepackage{fancyhdr}

% --- COLORES INSTITUCIONALES ---
\definecolor{valblue}{RGB}{0,112,192}
\definecolor{valdark}{RGB}{30,30,30}
\definecolor{valgray}{RGB}{240,240,240}

% --- FUENTE ---
\usepackage{helvet}
\renewcommand{\familydefault}{\sfdefault}

% --- CONFIGURACIÓN DE PÁGINA ---
\pagestyle{fancy}
\fancyhf{}
\fancyhead[L]{\footnotesize \color{gray} LICEO V.A.L. -- MATEMÁTICAS 6º}
\fancyhead[R]{\footnotesize \color{gray} CÓDIGO: 06M-18}
\fancyfoot[C]{\footnotesize \color{gray} Página \thepage}
\renewcommand{\headrulewidth}{0.4pt}
\renewcommand{\footrulewidth}{0pt}

% --- CONFIGURACIÓN DE RESPUESTAS CONDICIONALES ---
\newif\ifshowanswers
% CAMBIA A \showanswerstrue PARA VER EL SOLUCIONARIO, O \showanswersfalse PARA LA VERSIÓN DEL ESTUDIANTE
\showanswerstrue 

% Comando de respuestas corregido (recibe 1 solo parámetro)
\newcommand{\answer}[1]{
	\ifshowanswers
	\par\vspace{0.15cm}
	\begin{tcolorbox}[colback=blue!5!white, colframe=blue!60!white, arc=1mm, boxrule=0.5pt, title=\textbf{\footnotesize SOLUCIÓN (DOCENTE)}]
		\color{blue!80!black}\footnotesize #1
	\end{tcolorbox}
	\par\vspace{0.15cm}
	\else
	\par\vspace{0.3cm}
	\fi
}

\begin{document}
	
	\begin{center}
		{\Large \textbf{\color{valblue} LICEO V.A.L.}}\\[0.1cm]
		{\large \textbf{Taller de Repaso -- Trabajo en Cuaderno}}\\[0.1cm]
		{\small \textbf{Materia: Matemáticas 6º} \quad | \quad \textbf{Unidad: Estadística II y Probabilidad} \quad | \quad \textbf{Código: 06M-18}}
	\end{center}
	
	\vspace{0.2cm}
	\noindent \textbf{Nombre del Estudiante:} \underline{\hspace{8cm}} \hfill \textbf{Fecha:} \underline{\hspace{3cm}}
	
	\vspace{0.4cm}
	\begin{tcolorbox}[colback=valgray, colframe=valdark, arc=1mm, boxrule=0.8pt]
		\textit{Instrucciones: Transcribe y resuelve cada uno de los siguientes puntos en tu cuaderno de matemáticas de forma ordenada. Escribe los procedimientos completos, justificaciones y operaciones para validar tu aprendizaje.}
	\end{tcolorbox}
	
	\vspace{0.3cm}
	
	% --- PREGUNTA 1 ---
	\noindent \textbf{1. (Misión 1) [25\%]} Un grupo de estudiantes fue encuestado sobre el número de mascotas que tienen en sus hogares, obteniendo los siguientes datos:
	\begin{center}
		\textbf{1, \ 2, \ 0, \ 1, \ 3, \ 1, \ 2, \ 0, \ 1, \ 1, \ 0, \ 2}
	\end{center}
	\begin{enumerate}[label=\alph*)]
		\item Construye la tabla de frecuencias para este conjunto de datos.
		\item ¿Cuántos compañeros en total participaron en la encuesta? ¿Cuántos tienen exactamente 1 mascota?
	\end{enumerate}
	
	\answer{
		\textbf{a) Tabla de Frecuencias:}\\
		\begin{tabular}{cc}\toprule \textbf{Mascotas (Dato)} & \textbf{Frecuencia} \\\midrule 0 & 3 \\ 1 & 5 \\ 2 & 3 \\ 3 & 1 \\\midrule \textbf{Total} & \textbf{12} \\\bottomrule\end{tabular}\\[0.2cm]
		\textbf{b) Respuestas:} Participaron \textbf{12} compañeros (suma de frecuencias); exactamente \textbf{5} de ellos tienen 1 mascota.
	}
	
	% --- PREGUNTA 2 ---
	\noindent \textbf{2. (Misión 2) [25\%]} Utilizando la información recolectada en la Misión 1:
	\begin{enumerate}[label=\alph*)]
		\item Construye un diagrama de barras (recuerda nombrar los ejes).
		\item Identifica cuál es la \textbf{moda} de este conjunto de datos y explica qué significa ese valor en el contexto de la encuesta.
	\end{enumerate}
	
	\answer{
		\textbf{a)} El estudiante debe dibujar un gráfico donde el eje $X$ sea la "Cantidad de mascotas" (0, 1, 2, 3) y el eje $Y$ sea la "Frecuencia" (número de estudiantes, llegando al menos al 5). Las alturas deben ser 3, 5, 3 y 1 respectivamente. \\
		\textbf{b) Moda:} La moda es \textbf{1 mascota} (frecuencia de 5). Significa que lo más común en ese grupo de estudiantes es tener exactamente una mascota.
	}
	
	% --- PREGUNTA 3 ---
	\noindent \textbf{3. (Misión 3) [25\%]} Las notas obtenidas por Camila en sus evaluaciones de matemáticas durante el periodo fueron las siguientes: 
	\begin{center}
		\textbf{7, \ 8, \ 6, \ 9, \ 8, \ 10, \ 8}
	\end{center}
	\begin{enumerate}[label=\alph*)]
		\item Calcula la media aritmética (promedio) de sus notas. Escribe la operación completa.
		\item Encuentra la mediana de las notas de Camila. (Pista: ¡No olvides ordenar los datos primero!).
	\end{enumerate}
	
	\answer{
		\textbf{a) Media (Promedio):} Suma de los datos = $7 + 8 + 6 + 9 + 8 + 10 + 8 = 56$. Cantidad de datos ($n$) = 7. \\
		Media = $\frac{56}{7} = \textbf{8}$. \\
		\textbf{b) Mediana:} Ordenando los datos de menor a mayor: 6, 7, 8, \textbf{8}, 8, 9, 10. El dato central (posición 4) es el \textbf{8}. Por tanto, Mediana = 8.
	}
	
	% --- PREGUNTA 4 ---
	\noindent \textbf{4. (Misión 4) [25\%]} En una bolsa oscura de sorpresas hay \textbf{4 pelotas rojas, 3 pelotas azules y 5 pelotas verdes}. Si metes la mano y sacas una pelota al azar sin mirar:
	\begin{enumerate}[label=\alph*)]
		\item ¿Cuál es la probabilidad de sacar una pelota roja? (Exprésalo como fracción).
		\item ¿Qué color de pelota es menos probable que saques? ¿Por qué?
	\end{enumerate}
	
	\answer{
		\textbf{a) Probabilidad (Roja):} Casos favorables = 4 (rojas). Casos posibles (total de pelotas) = $4 + 3 + 5 = 12$. \\
		Probabilidad = $\mathbf{\frac{4}{12}}$ (Se puede simplificar a $\frac{1}{3}$).\\
		\textbf{b) Menos probable:} El color \textbf{azul} es el menos probable de sacar porque es el que tiene la menor cantidad (frecuencia) en la bolsa (solo 3 pelotas).
	}
	
\end{document}